% Registered A.47: full independent parent proof audit and code rerun PASS.
\subsection{Normalizers in every odd dihedral cover}

\begin{statement}{A.47}{Proved}{All-prime normalizer recognition for the odd dihedral family}
For every odd \(n\ge3\), every prime \(\ell\), and
\(X_n\) of A.43, the nontrivial closed procyclic
subgroups of \(\pi_1(X_{n,\bar K})^{(\ell)}\)
whose arithmetic normalizer has open Galois image are
exactly the nontrivial closed subgroups of conjugates
of the two cusp inertia groups. Together with A.43,
this proves both the full fixed-group and normalizer
assertions of Task A.5 for the entire odd dihedral
family, including composite \(n\).

The tame part uses the following general obstruction:
in compact-type reduction, a nonzero proper homology
class of an arithmetically normalized procyclic subgroup
cannot have nonzero projections to two different
positive-genus components, when the coefficient prime
differs from the residue characteristic.
\end{statement}

\begin{proof}
\paragraph{The compact-type component-support obstruction.}
We prove the stated general obstruction first. Work at
residue characteristic \(p\) with coefficient
\(\ell\ne p\), after finite extension so the
semistable model and its components are split. Write
\[
               V=\bigoplus_{i=1}^r V_i,
               \qquad\Omega=\sum_{i=1}^r\Omega_i,
                                                               \tag{A.47.1}
\]
for the proper homology and its local orthogonal
symplectic decomposition. The \(V_i\)'s belong to
the positive-genus components of the compact-type
special fibre. Choose an ordinary base point at a
smooth point on any root component. Root and orient
the graph away from it; let \(m_e>0\) be the edge
lengths and \(i_e\) the indicator that component
\(i\) is downstream from edge \(e\).

The actual based edge formula of
\cite[Definition 3.1.5 and Theorem 3.1.6]{fghBettsDogra}
gives, on the proper rational unipotent group, zero
first Johnson degree and first possible degree two
\[
 D(v_i)=[W_i,v_i],\qquad
 W_i=\sum_e m_ei_e\delta_e,\qquad
 \delta_e=\sum_j j_e\Omega_j.                        \tag{A.47.2}
\]
A nonzero tame inertia power multiplies this entire
derivation by a common nonzero scalar, which does not
affect the argument. The rational graded Lie algebra
through the required degrees is the usual surface
presentation
\[
                           \mathfrak l=\mathbb L(V)/(\Omega),
                                                               \tag{A.47.3}
\]
as used in A.25. The choice of representative of
outer monodromy changes it only by an inner term.

Suppose an arithmetically normalized procyclic
subgroup has nonzero proper homology class \(v\).
Its image in the proper pro-\(\ell\) group is still
nontrivial, with the same open arithmetic normalizer
image. Restrict the number field so that its absolute
Galois group lies in that image. The local inertia
action on \(V\) is trivial, because the proper
Jacobian has good reduction. Therefore any lift of
inertia in the normalizer acts on the procyclic
subgroup by exponent one: its exponent is also the
scalar on the nonzero homology vector \(v\).
An actual nonzero power of tame inertia must thus
carry a generator \(u\) to a conjugate of itself
under any representative of its proper outer action.

We recall why this imposes a precise degree-three
condition without any integral lifting assumption.
Use the representative (A.47.2), with zero Johnson
degree one. Comparing degree two in the rational
conjugacy equation gives
\[
                 [q_1,v]=0\quad\text{in }\Lambda^2V/\mathbf Q_\ell\Omega,
                                                               \tag{A.47.4}
\]
where \(q_1\) is the first term of a conjugator.
If two positive-genus components occur, the total
genus is at least two. The decomposable bivector
\(q_1\wedge v\), of rank at most two, cannot be
a nonzero multiple of the nondegenerate rank-
\(2g\) bivector \(\Omega\). Thus \(q_1\) is
proportional to \(v\). Multiplying the rational
conjugator on the right by the corresponding inverse
rational power of \(u\) removes \(q_1\) without
changing its conjugation of \(u\). In degree three,
all terms involving higher parts of \(\log u\)
now disappear, and one obtains the necessary condition
\[
                              D(v)\in[v,\mathfrak l_2].
                                                               \tag{A.47.5}
\]
This uses a rational quotient only for a necessary
equation. It claims neither that a rational power
lies in the integral group nor that an actual
normalizer extension splits. Adding an inner
representative to \(D\) does not change (A.47.5).

Suppose \(v_i\ne0\) and \(v_j\ne0\) for distinct
components \(i,j\). Extend each of these vectors
to a symplectic basis of its block. There is a Lie
map from (A.47.3) to the free Lie algebra on \(x,y\)
which sends
\[
 v_i\longmapsto x,\quad v_j\longmapsto-x,\qquad
 \Omega_i\longmapsto[x,y],\quad
 \Omega_j\longmapsto-[x,y],                         \tag{A.47.6}
\]
and all other blocks to zero. Within either selected
block choose the companion of its indicated vector
so that the first symplectic pair maps to the shown
commutator, and kill all its other pairs. This
constructs the map even if a component has genus
greater than one. It kills both \(v\) and the
surface relation \(\Omega\).

Put \(Z=[x,y]\). Formula (A.47.2) now gives
\[
 \begin{split}
 D(v)&\longmapsto
          \left(\sum_e m_e(i_e-j_e)^2\right)[Z,x]\\
      &=\operatorname{dist}(C_i,C_j)\,[[x,y],x]\ne0.
 \end{split}                                                   \tag{A.47.7}
\]
An edge has \((i_e-j_e)^2=1\) exactly when it
separates the two components, so the coefficient is
the positive distance between distinct vertices.
The Lie word is nonzero: in the free associative
algebra it is \(2xyx-yxx-xxy\). But (A.47.6)
kills the right side of (A.47.5), since it kills
\(v\). This contradiction proves that a candidate
nonzero homology class is supported on a single
positive-genus component. The proof works for every
compact-type tree, not only for a star.

\paragraph{Apply this obstruction to the odd dihedral covers.}
Suppose that a noncuspidal nontrivial procyclic subgroup
of the geometric pro-\(\ell\) group of \(X_n\)
has open arithmetic normalizer. By A.11 a generator
has nonzero proper homology class
\[
                         v\in V_\ell\operatorname{Jac}(C_n).
\]
The actual punctured morphism
\(f=[1+i]\circ q_n:X_n\to U\) of A.43 implies
\(f_*v=0\). Indeed, if that proper image were
nonzero, the image of the procyclic subgroup in the
target pro-\(\ell\) group would be nontrivial and
would still have open arithmetic normalizer. A.12
puts it in target cusp inertia, whose proper homology
is zero, a contradiction. This step uses the actual
map of punctured curves and the induced arithmetic
homomorphism, not just a homomorphism of Jacobians.

Let \(\ell\) be odd and use the actual model at
two from A.43. Its \(n\) positive-genus components
are copies of the original supersingular elliptic
special fibre. The normalization proof identifies
the map of each such component to the main component
of the base model with the identity. Hence, under
compact-type specialization,
\[
             V_\ell\operatorname{Jac}(C_n)
                         =\bigoplus_{i=1}^n V_i,
\]
and every restriction
\(f_*|_{V_i}:V_i\to V_\ell E\) is an isomorphism:
it is the component identity followed by the
isogeny \([1+i]\). This compatibility also follows
by functoriality of the good Jacobian models. The
possible inseparability of that special isogeny at
two does not affect its prime-to-two rational Tate
module isomorphism.

A nonzero vector in \(\ker f_*\) cannot be
supported on one \(V_i\), since that restriction
is injective. Thus \(v\) has nonzero projections
onto at least two components. The compact-type
component-support obstruction just proved, with
\(p=2\), contradicts this. It excludes every
noncuspidal normalized procyclic subgroup at every
odd coefficient prime, including the primes dividing
\(n\); only nonzero rational coefficients were used.

At \(\ell=2\), A.43 proves that the proper Jacobian
has potentially good reduction at two with special
two-rank zero. The general normalizer criterion
A.46 therefore excludes a noncuspidal pro-two
subgroup directly. It applies over every finite
field extension needed to contain the arithmetic
normalizer image.

Conversely, after making a cusp rational its inertia
has a decomposition group with open arithmetic
projection. Every nonzero closed subgroup of its
pro-\(\ell\) inertia is characteristic there and
has the same open normalizer property. This proves
the exact classification in the statement.

The program \texttt{code/a\_odd\_dihedral\_normalizers.py}
checks the distance coefficient on every labelled
tree of small order, the exact two-block Lie witness,
and the component-map kernel implication for the
explicit stars. The obstruction above is uniform
in the number of components and in the tree.
\end{proof}
